\section{Bounding Hochschild homology of \texorpdfstring{$\BP\langle n\rangle$}{BP n }}\label{sec:bounding}
The goal of this section is to prove the following result.  
\begin{proposition}\label{prop:brunss}
There is a multiplicative Brun spectral sequence 
\begin{equation}\label{BrunSS}
\THH_*(\BP\langle n-1 \rangle)\langle \sigma v_{n}\rangle \implies \THH_*(\BP\langle n\rangle,\BP\langle n-1\rangle)\,.
\end{equation}
which is a spectral sequence of $(\BP\langle n-1 \rangle\otimes_{\BP\langle n\rangle}\BP\langle n-1 \rangle)_*$-algebras and the differentials are $(p,v_1,\cdots ,v_{n-1})$-linear and $\sigma v_{n}$-linear. The $d_1$-differential is given by the cap product of a certain class in $\THC^{2p^{n+1}}(\BP\langle n-1\rangle)$.  
\end{proposition}

To prove this, we first need the following lemma. 
\begin{lemma}\label{lem:BPn-rel-cooperations}
Suppose $-1\le m\le n$. We compute that 
\[\pi_*\BP\langle m \rangle\otimes_{\BP\langle n\rangle}\BP\langle m \rangle=\mathbb{Z}_{(p)}[v_1,\cdots ,v_{m}]\langle \sigma v_{m+1} , \cdots ,\sigma v_{n} \rangle \,. \]
when $m \geq 0$ and 
\[
\pi_*\BP\langle -1 \rangle\otimes_{\BP\langle n\rangle}\BP\langle -1\rangle=\mathbb{F}_p\langle \sigma v_{m+1} , \cdots ,\sigma v_{n} \rangle \,. 
\]
\end{lemma}
\begin{proof}
This follows directly from the K\"unneth spectral sequence
\[ 
\Tor_*^{\BP\langle n\rangle_*}(\pi_*\BP\langle m\rangle,\pi_*\BP\langle m\rangle)\Longrightarrow \pi_*\BP\langle m\rangle\otimes_{\BP\langle n\rangle}\BP\langle m\rangle
\]
which collapses at the $\mathrm{E}_2$-page because all differentials on algebra generators land in trivial bidegrees and $\mathbb{Z}_{(p)}[v_1,\cdots ,v_{m}]$ splits off. 
\end{proof}

By the same argument as \cite{AHL10}*{Proposition~5.2}, we compute that 
\[ \pi_*H\mathbb{F}_p\otimes_{\BP}\BP\langle n\rangle=\mathbb{F}_p\langle \overline{\tau}_{n+1},\overline{\tau}_{n+2},\cdots \rangle\]
and using the universal coefficient spectral sequence we deduce that 
\[ \THC^*(\BP\langle n\rangle/\BP ; \mathbb{F}_p)=\mathbb{F}_p[e_{n+1},e_{n+2},\cdots ]\]
where $\varepsilon_i$ corresponds to $\tau_i$. Since these classes are in even degrees, the Bockstein spectral sequences collapse leading to
\[ 
\THC^*(\BP\langle n\rangle/\BP )=\BP\langle n\rangle_*[[e_{n+1},e_{n+2},\cdots]]
\]
The computation that 
\[ 
\THC^*(\BP\langle n\rangle ; \mathbb{F}_p)=\mathbb{F}_p\langle \lambda_1^{\vee},\cdots ,\lambda_{n+1}^{\vee})\otimes \Gamma ((\mu^{p^{n+1}})^{\vee})
\]
is immediate from~\cite{ACH24}*{Proposition~2.9}. By examination of the map of universal coefficient spectral sequences, we can also deduce that the map
\[ \THC^*(\BP\langle n\rangle/\BP)\to \THC^*(\BP\langle n\rangle)\]
sends $e_{n+1}$ to $(\mu^{p^{n+1}})^{\vee}$. and consequently, there exists a lift $e_{n+1}\in \THC^{2p^{n+1}}(\BP\langle n\rangle)$ of $(\mu^{p^{n+1}})^{\vee}$ by the commuting diagram
\[
\begin{tikzcd}
\THC^*(\BP\langle n\rangle/\BP) \arrow{d} \arrow{r}& \THC^*(\BP\langle n\rangle/\BP;\mathbb{F}_p) \arrow{d}\\
\THC^*(\BP\langle n\rangle)\arrow{r} & \THC^*(\BP\langle n\rangle;\mathbb{F}_p) \,.
\end{tikzcd}
\]

\begin{proof}[Proof of Proposition~\ref{prop:brunss}]
Let $F$ denote the fiber of the map 
\[ 
\BP\langle n-1\rangle\otimes_{\BP\langle n\rangle}\BP\langle n-1\rangle \longrightarrow \BP\langle n-1\rangle \otimes_{\BP\langle n-1\rangle}\BP\langle n-1\rangle =\BP\langle n-1\rangle 
\]
of right $\mathrm{BP}\langle n-1\rangle \otimes \mathrm{BP}\langle n-1\rangle^{\op}$-modules and left $\BP\langle n-1\rangle\otimes_{\BP\langle n\rangle}\BP\langle n-1\rangle$-modules. The right $\mathrm{BP}\langle n-1\rangle \otimes \mathrm{BP}\langle n-1\rangle^{\op}$-module structure can also be considered as a $\mathrm{BP}\langle n-1\rangle$-bimodule structure. As a left $\mathrm{BP}\langle n-1\rangle$-module $F\simeq \mathrm{BP}\langle n-1\rangle$. To see this, note that $\mathrm{BP}\langle n-1 \rangle=\mathrm{BP}\langle n\rangle/v_n$ and therefore 
\[ \mathrm{BP}\langle n-1 \rangle\otimes_{\mathrm{BP}\langle n\rangle} \mathrm{BP}\langle n\rangle/v_n=\mathrm{BP}\langle n-1 \rangle/v_n\]
but $v_n$ acts trivially on $\mathrm{BP}\langle n-1 \rangle$ so $\mathrm{BP}\langle n-1 \rangle/v_n\simeq \mathrm{BP}\langle n-1 \rangle\vee \Sigma^{2p^{n-1}}\mathrm{BP}\langle n-1 \rangle$ as a left $\mathrm{BP}\langle n-1 \rangle$-module so the fiber can be identified with $\mathrm{BP}\langle n-1 \rangle$ as a left $\mathrm{BP}\langle n-1 \rangle$-module. A similar argument with right 
$\mathrm{BP}\langle n\rangle$-module structures implies that the fiber is equivalent to $\mathrm{BP}\langle n-1\rangle$ as a right $\mathrm{BP}\langle n-1\rangle$-module and consequently as a $\mathrm{BP}\langle n-1\rangle\otimes \mathrm{BP}\langle n-1\rangle^{\op}$-module. Note that we are not claiming that $\mathrm{BP}\langle n-1\rangle \otimes_{\mathrm{BP}\langle n\rangle}\mathrm{BP}\langle n-1\rangle$ splits as a $\mathrm{BP}\langle n-1\rangle \otimes_{\mathrm{BP}\langle n\rangle}\mathrm{BP}\langle n-1\rangle^{\op}$ bimodule. 

However, we do produce a fiber sequence 
\[ 
\Sigma^{2p^{n}-1}\THH(\BP\langle n-1 \rangle)\to  \THH( \mathrm{BP}\langle n \rangle;\mathrm{BP}\langle n-1 \rangle)\longrightarrow \THH(\BP\langle n-1 \rangle)
\]
of left $\BP\langle n-1\rangle \otimes_{\BP\langle n \rangle }\BP\langle n-1\rangle $-modules 
and the long exact sequence associated to this fiber sequence may be regarded as spectral sequence, which we call the Brun spectral sequence following~\cite{Hon20}. In fact, since the map 
\[ \mathrm{BP}\langle n-1\rangle \to \Sigma^{2p^{n}}\mathrm{BP}\langle n-1\rangle 
\]
is a $\mathrm{BP}\langle n-1\rangle \otimes \mathrm{BP}\langle n-1\rangle^{\op}$-module map, it is an element in topological Hochschild cohomology $\mathrm{THC}(\mathrm{BP}\langle n-1\rangle)$ and therefore it is a $\mathbb{E}_1$-derivation so the associated spectral sequence is multiplicative. The $d_1$-differential is therefore given by capping with this class, which can be identified with a choice of lift $\varepsilon_{n+1}$ of the class $(\mu^{p^{n+1}})^{\vee}$. 
\end{proof}
\begin{remark}
Given the equivalence  
\[ \THH(A;B)\simeq \THH(B;B\otimes_{A}B)\]
one can filter $B\otimes_{A}B$ by its Whitehead filtration $\tau_{\ge \bullet}$ to produce a spectral sequence. Such a spectral sequence was originally considered by Brun~\cite{Bru00} in the setting of FSP's and was extended to the setting of commutative ring spectra by H\"oning~\cite{Hon20}, for applications see also \cite{Cha25} and~\cite{Hys23}. Although we consider a more general filtration of $B\otimes_{A}B$ above, we were inspired by the work of Brun and H\"oning and therefore we still call this a Brun spectral sequence. 
\end{remark}
\begin{example}
Consider the case $n=0$. Then there is a spectral sequence 
\[ 
\THH_{*}(\mathbb{F}_{p})\otimes \Lambda (\sigma v_{0})\implies \THH_{*}(\mathbb{Z}_{(p)};\mathbb{F}_{p}) \,.
\]
This spectral sequence has a differential $d_{1}(\mu)=\sigma v_{0}$ and no further differentials except those generated by the Leibniz rule yielding the known answer $\mathbb{F}_p[\mu^p]\langle \sigma v_0 \mu^{p-1}\rangle$. 
Already, this shows that although $\mathbb{F}_p\otimes_{\mathbb{Z}_{(p)}}\mathbb{F}_p$ splits as an $\mathbb{F}_p$-module it does not split as a $\mathbb{F}_p\otimes \mathbb{F}_p$-module. This is perhaps not surprising because $\pi_*\mathbb{F}_p\otimes_{\mathbb{Z}_{(p)}}\mathbb{F}_p=\Lambda (\sigma v_0)$ and the $k$-invariant $\mathbb{F}_p\to \Sigma^{2}\mathbb{F}_p$ is the non-trivial class in degree $2$ in topological Hochschild cohomology of $\mathbb{F}_p$ which is dual to $\mu$ (up to a unit).
\end{example}

\begin{example}
Consider the case $n=1$. Then there is a spectral sequence 
\[ 
\THH_{*}(\mathbb{Z}_{(p)})\otimes \Lambda (\sigma v_{1})\implies \THH_{*}(\ell;\mathbb{Z}_{p}) \,.
\]
We know $\sigma v_1$ survives because it generates a torsion free class and $\THH_*(\mathbb{Z})$ is torsion above degree zero. Specifically, 
\[
\THH_*(\mathbb{Z}_{(p)})=\begin{cases} \mathbb{Z}_{(p)} & \text{ if } n=0 \\ 
\mathbb{Z}/p^{\nu_p(k)+1}\{\lambda_1\mu_1^{k-1}\} & \text{ if } n=2pk-1>0 \\ 
0 & \text{ otherwise.}
\end{cases}
\]
By comparing to~\cite{AHL10} for example, we can determine differentials 
\[
d_1(\lambda_1\mu^{k})\dot{=}p^{\nu_p(k)}\lambda_1\mu^{k-1}\sigma v_1 
\]
for all integers $k$. 
In particular, when $k=ip-1$ for some integer $i\ge 0$ then 
\[ 
\ker
(d_1 :\mathbb{Z}/p^{\nu_p(i)+2}\{\lambda_1\mu_1^{ip-1}\}\to \mathbb{Z}/p\{\lambda_1\mu_1^{ip-2}\sigma v_1) = \mathbb{Z}/p^{\nu_p(i)+1}\{\lambda_1\mu^{ip-1}\} \,.
\]
and when $k=pi$ for some integer $i\ge 0$ then 
\[ 
\textup{coker} 
(d_1 :\mathbb{Z}/p\{\lambda_1\mu_1^{ip}\}\to \mathbb{Z}/p^{\nu_p(i)+2}\{\lambda_1 \mu_1^{ip-1}\sigma v_1 \}) = \mathbb{Z}/p^{\nu_p(i)+1}\{\sigma v_1\lambda_1\mu_1^{ip-1}\}
\]

This accounts for the classes $a_i=p\lambda_1\mu_1^{ip-1}$ and $b_i=\sigma v_1\lambda_1\mu_1^{ip-1}$ respectively. Otherwise the kernel and cokernel are trivial. In this case, the $k$-invariant is a map 
\[ H\mathbb{Z}\to \Sigma^{2p} H\mathbb{Z}\]
of $ H\mathbb{Z}\otimes H\mathbb{Z}^{\op}$-modules, corresponding to a generator of $\mathrm{THC}^{2p}(\bZ)=\mathbb{Z}/p$. Here we use the equivalence 
\[ 
\map_{H\bZ\otimes H\bZ^{\op}}(H\bZ,H\bZ)\simeq \map_{H\bZ}(\THH(\bZ),H\bZ) \,,
\]
the universal coefficient spectral sequence 
\[ 
\Ext_{\mathbb{Z}}^{*,*}(\THH_*(\mathbb{Z}),\mathbb{Z})\implies \THC^*(\mathbb{Z})
\]
associated to the equivalence, and the known computation 
\[
    \THH_n(\mathbb{Z})=\begin{cases} \mathbb{Z} &n=0 \\ 
    \mathbb{Z}/(k-1) &n=2k-1 \\ 
    0 & \text{otherwise}
    \end{cases}
\]
to produce the identity $\mathrm{THC}^{2p}(\mathbb{Z})=\mathbb{Z}/p$
\end{example}
